\subsection{例子：一维的情形}

在本节中，考虑一个一维的 Noether 整环 A，它具有对偶化模 \(\Omega\)。用 K 表示 A 的分式域，用 V 表示 K-向量空间 \(K \otimes_{A} \Omega\)。

典范同态 \(\Omega\to V\) 是单射的，且 K-向量空间 \(V\) 的维数为 1（\(n^{\circ}\) 2, 命题 4 的推论 1）；把 \(\Omega\) 与 \(V\) 的一个子 A-模等同起来。

命题 8. A-模 V/Ω 是 Matlis 模。

考虑正合列

\[
0 \to \Omega \to \mathrm{V} \to \mathrm{V} / \Omega \to 0;
\]

A-模 V 是内射的（A, X, 第 18 页，例 1），且有 \(\mathrm{di}_{\mathrm{A}}(\Omega) = 1\)（\(n^{\circ} 1\)，命题 1）。由此一方面推出 \(V/\Omega\) 是内射的（\(\S 3\)，\(n^{\circ} 1\)，命题 1），另一方面推出：对 A 的每个极大理想 \(\mathfrak{m}\)，\(\Lambda/\mathfrak{m}\)-向量空间 \(\mathrm{Hom}_{\mathrm{A}}(\mathrm{A}/\mathfrak{m}, \mathrm{V}/\Omega)\) 同构于 \(\mathrm{Ext}_{\mathrm{A}}^{1}(\mathrm{A}/\mathfrak{m}, \Omega)\)，故维数为 1。由于 \(V/\Omega\) 是扭模，它的伴随素理想都是极大的；这就证明了命题（\(\S 8\)，\(n^{\circ} 4\)）。

设 M 为 A-模；依照前引文，用 D(M) 表示 A-模 \(\mathrm{Hom}_{\mathrm{A}}(\mathrm{M},\mathrm{V}/\Omega)\)。可以把第 5 节的构造应用于此，取 I 为复形

\[
\dots 0 \longrightarrow \mathrm{V} \xrightarrow {p} \mathrm{V} / \Omega \longrightarrow 0 \dots
\]

其中 V 置于次数 0，而 p 表示典范满射。复形 D(M) 为

\[
\dots 0 \longrightarrow \operatorname{Hom} _ {\mathrm{A}} (\mathrm{M}, \mathrm{V}) \xrightarrow {p _ {\mathrm{M}}} \mathrm{D} (\mathrm{M}) \longrightarrow 0 \dots ,
\]

其中 \(p_{\mathbf{M}} = \mathrm{Hom}(\mathbf{1}_{\mathbf{M}}, p)\)。我们有分次 A-模的一个典范同构 \(\varphi(\mathbf{M}, \mathbf{I}) : \mathrm{H}(\mathbf{D}(\mathbf{M})) \longrightarrow \mathrm{Ext}_{\mathbf{A}}(\mathbf{M}, \Omega)\)（A, X, 第 100 页，定理 1）。

当 M 是扭模时，模 \(\mathbf{D}(\mathbf{M})^{0} = \operatorname{Hom}_{\mathbf{A}}(\mathbf{M}, \mathbf{V})\) 为零，且 \(\varphi(\mathbf{M}, \mathbf{I})\) 是从 D(M) 到 \(\operatorname{Ext}_{\Lambda}^{1}(\mathbf{M}, \Omega)\) 的同构；态射 \(\alpha_{M}: M \to \mathbf{D}(\mathbf{D}(\mathbf{M}))\) 不是别的，正是 A-模的典范同态

\[
\alpha_ {\mathrm{M}}: \mathrm{M} \longrightarrow \mathrm{D} (\mathrm{D(M)})
\]

它定义于 § 8, 第 3 节，且当 M 是有限型的（即有限长度的）时是同构。在此情形我们重新得到前引文中的状况。

回到一般情形，并设 A-模 M 是有限型的。于是 D(M) 是扭模，故 A-模 \(\mathbf{D}(\mathbf{D}(\mathbf{M}))^{-1} = \mathrm{Hom}_{\mathrm{A}}(\mathrm{D}(\mathrm{M}), \mathrm{V})\) 为零。另一方面，A-模 \(\mathrm{Hom}_{\mathrm{A}}(\mathbf{D}(\mathbf{M})^{0}, \mathrm{V}) = \mathrm{Hom}_{\mathrm{A}}(\mathrm{Hom}_{\mathrm{A}}(\mathrm{M}, \mathrm{V}), \mathrm{V})\) 自然等同于 \(K \otimes_{A} M\)，于是同态

\[
\operatorname{Hom} (1, p): \operatorname{Hom} _ {\mathrm{A}} (\operatorname{Hom} _ {\mathrm{A}} (\mathrm{M}, \mathrm{V}), \mathrm{V}) \longrightarrow \mathrm{D} (\operatorname{Hom} _ {\mathrm{A}} (\mathrm{M}, \mathrm{V}))
\]

等同于映射

\[
j: \mathrm{K} \otimes_ {\mathrm{A}} \mathrm{M} \longrightarrow \mathrm{D} (\operatorname{Hom} _ {\mathrm{A}} (\mathrm{M}, \mathrm{V}))
\]

它满足 \(j(\lambda \otimes m)(f) = p(\lambda f(m))\)（对 \(\lambda \in \mathbf{K}\)，\(m \in \mathbf{M}\)，\(f \in \mathrm{Hom}_{\mathrm{A}}(\mathrm{M},\mathrm{V})\)）。于是第 5 节的定理 1 就转化为序列的正合性

\[
0 \longrightarrow \mathrm{M} \xrightarrow {(i , \alpha_ {\mathrm{M}})} (\mathrm{K} \otimes_ {\mathrm{A}} \mathrm{M}) \oplus \mathrm{D} (\mathrm{D} (\mathrm{M})) \xrightarrow {(j , - \mathrm{D} (p _ {\mathrm{M}}))} \mathrm{D} (\operatorname{Hom} _ {\Lambda} (\mathrm{M}, \mathrm{V})) \longrightarrow 0,
\]

其中 \(i\) 表示从 M 到 \(\mathrm{K} \otimes_{\mathrm{A}} \mathrm{M}\) 的典范映射。\(i\) 的核等同于 M 的扭子模 \(\mathrm{T(M)}\)，它的余核等同于 \((\mathrm{K/A}) \otimes_{\mathrm{A}} \mathrm{M}\)。考虑行正合的交换图表

\[
\begin{array}{c c c c c c c} 0 \to \mathrm{T(M)} & \longrightarrow & \mathrm{M} & \stackrel {{\imath}} {{\longrightarrow}} & \mathrm {K\otimes_ {A} M} & \longrightarrow & (\mathrm{K/A}) \otimes_ {\mathrm{A}} \mathrm{M} \longrightarrow 0 \\ & & ^ {\alpha_ {\mathrm{M}}} \Bigg \downarrow & & \Bigg \downarrow_ {j} \\ 0 \to \mathrm{D(Ext} _ {\mathrm{A}} ^ {1} (\mathrm{M}, \Omega)) & \longrightarrow & \mathrm{D(D(M))} & \stackrel {{\mathrm{D} (p _ {\mathrm{M}})}} {{\longrightarrow}} & \mathrm{D(Hom} _ {\mathrm{A}} (\mathrm{M}, \mathrm{V})) & \longrightarrow & \mathrm{D(Hom} _ {\mathrm{A}} (\mathrm{M}, \Omega)) \longrightarrow 0 \end{array}
\]

其中第二行是由正合列经 Matlis 对偶得到的

\[
0 \rightarrow \operatorname{Hom} _ {\mathrm{A}} (\mathrm{M}, \Omega) \longrightarrow \operatorname{Hom} _ {\mathrm{A}} (\mathrm{M}, \mathrm{V}) \xrightarrow {p _ {\mathrm{M}}} \operatorname{Hom} _ {\mathrm{A}} (\mathrm{M}, \mathrm{V} / \Omega) \longrightarrow \operatorname{Ext} _ {\Lambda} ^ {1} (\mathrm{M}, \Omega) \rightarrow 0.
\]

于是定理 1 的含义是：A-模的同态

\[
\gamma^ {0} (\mathrm{M}): \mathrm{T} (\mathrm{M}) \longrightarrow \mathrm{D} \left(\operatorname{Ext} _ {\mathrm{A}} ^ {1} (\mathrm{M}, \boldsymbol {\Omega})\right) e t \quad \gamma^ {1} (\mathrm{M}): (\mathrm{K} / \mathrm{A}) \otimes_ {\Lambda} \mathrm{M} \longrightarrow \mathrm{D} \left(\operatorname{Hom} _ {\mathrm{A}} (\mathrm{M}, \boldsymbol {\Omega})\right)
\]

——它们分别由 \(\alpha_{\mathrm{M}}\) 与 \(j\) 导出——是双射的。由于 A-模 T(M) 是有限长度的，A-模 \(\operatorname{Ext}_{\mathrm{A}}^{1}(\mathrm{M},\Omega)\) 是有限长度的，且等同于 Matlis 对偶 D(T(M))，并且在群 Z0(A) 中有 \([\operatorname{Ext}_{\mathrm{A}}^{1}(\mathrm{M},\Omega)] = [\operatorname{T}(\mathrm{M})]\) 以及 \(\operatorname{long}_{\mathrm{A}}(\operatorname{Ext}_{\mathrm{A}}^{1}(\mathrm{M},\Omega)) = \operatorname{long}_{\mathrm{A}}(\operatorname{T}(\mathrm{M}))\)（\(\S 8\)，第 4 节，命题 4）。另一方面，取 M = A，便得到典范同构 \(\gamma^{1}(\mathrm{A}):K / A\to D(\Omega)\)。

设 B 为 K 的一个包含 A 的子环，且在 \(\Lambda\) 上是有限的。对 A 的每个极大理想 \(\mathfrak{m}\)，有 \(\mathrm{prof}_{\mathrm{A}_{\mathfrak{m}}}(B_{\mathfrak{m}}) = \dim_{A_{\mathfrak{m}}}(B_{\mathfrak{m}}) = 1\)（\(\S 1\)，\(n^{\circ} 1\)，注 2 及 VIII, \(\S 2\)，\(n^{\circ} 3\)，定理 1），故 B 是 Macaulay A-模。于是，B-模 \(\Omega_{\mathrm{B}} = \mathrm{Hom}_{\mathrm{A}}(B, \Omega)\) 是对偶化的（\(n^{\circ} 3\)，注）。从 \(\Omega_{\mathrm{B}} = \mathrm{Hom}_{\mathrm{A}}(B, \Omega)\) 到 \(\Omega = \mathrm{Hom}_{\mathrm{A}}(A, \Omega)\) 的典范映射是单射的；它的像由 \(\omega\) 的这样的元素 \(\Omega\) 组成：V 的子 B-模 B\(\omega\) 含于 \(\Omega\)。于是 \(\Omega_{\mathrm{B}}\) 等同于 \(\Omega\) 的最大子 B-模。A-模 B/\(\Lambda\) 是有限长度的；正合列

\[
0 \to \Omega_ {\mathrm{B}} \to \Omega \to \operatorname{Ext} _ {\mathrm{A}} ^ {1} (\mathrm{B/A}, \Omega) \to 0
\]

使我们可以把 \(\Omega/\Omega_{\mathrm{B}}\) 与 \(\operatorname{Ext}_{\mathrm{A}}^{1}(\mathrm{B}/\mathrm{A},\Omega)\) 等同，从而由前所述与 \(\mathrm{D}(\mathrm{B}/\mathrm{A})\) 等同。特别地，在 \([\mathrm{B}/\mathrm{A}] = [\Omega/\Omega_{\mathrm{B}}]\) 中有 \(Z_{0}(\mathrm{A})\) 以及 \(\operatorname{Ann}_{\mathrm{A}}(\mathrm{B}/\mathrm{A}) = \operatorname{Ann}_{\mathrm{A}}(\Omega/\Omega_{\mathrm{B}})\)（\(\S 8\)，\(\mathfrak{n}^{\circ}4\)，命题 4）。

理想 \(\mathfrak{c} = \mathrm{Ann}_{\mathrm{A}}(\mathrm{B / A})\) 是传递子 A : B，即（VII, § 1, 第 1 节）K 的元素 \(x\)（满足 \(xB \subset \Lambda\)）的集合。它是

A 的（非零）理想，也是 B 的理想；事实上它是 B 的含于 A 的最大理想。由于 \(\Omega_{\mathrm{B}} : \Omega \subset \Omega : \Omega = \mathrm{A} (\mathrm{n}^{\circ} 4, \text{prop. } 7, \mathrm{b}))\)，有 Ann \(_{\mathrm{A}}(\Omega/\Omega_{\mathrm{B}}) = \Omega_{\mathrm{B}} : \Omega\)，最终得到

\[
\mathfrak {c} = \operatorname{Ann} _ {\mathrm{A}} (\mathrm{B/A}) = \operatorname{Ann} _ {\mathrm{A}} (\Omega / \Omega_ {\mathrm{B}}) = \Omega_ {\mathrm{B}}: \Omega .
\]

由于 \(\Omega_{\mathrm{B}}\) 是 B-模，关系 \(x\Omega \subset \Omega_{\mathrm{B}}\) 等价于 \(xB\Omega \subset \Omega_{\mathrm{B}}\)，故又有 \(\mathfrak{c} = \Omega_{\mathrm{B}}: \mathrm{B}\Omega\)。

现在我们把前面的讨论特殊化到 B 是 A 的整闭包的情形；当环 A 是 Japanese 环（IX, § 4, 第 1 节，定义 1）时，B 是有限生成 A-模这一假设得到满足，而当 A 是局部且完备的环（前引文，第 2 节，定理 2），或基本上是在域上有限型的（前引文，第 1 节，注 2 及例）时就属于这种情形。此时环 B 是 Dedekind 环（VII, § 2, 第 2 节，定理 1），且无扭的 B-模 \(\Omega_{\mathrm{B}}\)、B \(\Omega\) 与 c 都是秩为 1 的投射模（VII, § 4, 第 10 节，命题 22）。于是关系 \(c = \Omega_{\mathrm{B}} : \mathrm{B}\Omega\) 的含义是：由 K 在 V 上的作用导出的线性映射 \(c \otimes_{\mathrm{B}} \mathrm{B}\Omega \to \Omega_{\mathrm{B}}\) 是同构（II, § 5, 第 6 节，命题 11）。特别地，有 \(\Omega_{\mathrm{B}} = c(\mathrm{B}\Omega) = c\Omega\)。

命题 9. 设 B 为 A 的整闭包，并令 \(\mathfrak{c} = \mathrm{A}:\mathrm{B}\)。假设 B 是有限生成的 A-模。在 \([\mathrm{B} / \mathrm{c}] \leqslant 2[\mathrm{B} / \mathrm{A}]\) 中有不等式 \(\mathrm{Z}_0(\mathrm{A})\)。为使等号成立，必须且只须 A 是 Gorenstein 环。

我们有 \([B/c] = [B/A] + [A/c]\)，故所考虑的不等式等价于 \([A/c] \leqslant [B/A]\)。

\begin{enumerate}
\def\labelenumi{\Alph{enumi})}
\item
  对 A 的每个极大理想 m，\(A_{m}\) 的整闭包是 \(B_{m}\)（V, § 1, 第 5 节，推论 1），且有 \(c_{m} = A_{m} : B_{m}\)。此外，按定义有 \([B/c] = \sum_{m} \text{long}_{A_{m}}(B_{m}/c_{m})[m]\) 与 \([B/A] = \sum_{m} \text{long}_{A_{m}}(B_{m}/A_{m})[m]\)。这就化归为环 A 是局部环的情形来证明命题，以下即作此假设。在此情形，有序群 \(Z_{0}(A)\) 典范地等同于 Z，于是有限长度模的类就是它的长度。环 B 是半局部的，且 B-模 BΩ 是秩为 1 的自由模（II, § 5, 第 3 节，命题 5）。
\item
  若 A 是 Gorenstein 环，则 A-模 \(\Omega\) 是秩为 1 的自由模（第 4 节，命题 7 的推论 1）；选取 \(\omega\) 的一个生成元 \(\Omega\)。有 \(\Omega = A\omega\) 与 \(\Omega_{\mathrm{B}} = c\Omega = c\omega\)，从而 \(\text{long}(A/c) = \text{long}(\Omega/\Omega_{\mathrm{B}}) = \text{long}(B/A)\)。
\item
  设剩余域 \(\kappa_{\mathrm{A}}\) 是无限的。对 B 的每个极大理想 \(\mathfrak{n}\)，用 L(n) 表示 B/n-向量空间（一维的）BΩ/nBΩ，并用 pr\(_{\mathfrak{n}}\) 表示从 \(\bigoplus_{\mathfrak{n}}\) L(n) 到 L(n) 上的典范投影。设 \(\varphi : \Omega \longrightarrow \bigoplus_{\mathfrak{n}}\) L(n) 是 \(\Omega\) 经典范同态 BΩ \(\longrightarrow \bigoplus_{\mathfrak{n}}\) L(n) 而得到的映射。\(\varphi\) 的像是 \(\kappa_{\mathrm{A}}\)-向量空间 \(\bigoplus_{\mathfrak{n}}\) L(n) 的一个子空间；它不含于 Ker pr \(_{\mathfrak{n}}\)，因为否则将有 \(\Omega \subset \mathfrak{n}\mathrm{B}\Omega\)，从而 BΩ \(\subset \mathfrak{n}\mathrm{B}\Omega\)，这是矛盾的。于是 \(\varphi\) 的像不含于诸 Ker pr \(_{\mathfrak{n}}\) 的并（A, V, 第 40 页，引理 1）；因此存在一个元素 \(\omega\)（属于 \(\Omega\)），它在 BΩ/nBΩ 中的像对每个 n 都非零，这蕴含 \(\omega\) 生成 B-模 BΩ（II, § 3, 第 3 节，命题 11）。
\end{enumerate}

设 \(a \in \mathbf{A}\)；若 \(a\omega\) 属于 \(\Omega_{\mathrm{B}}\)，则有 \(a\mathrm{B}\omega \subset \Omega_{\mathrm{B}}\)，从而 \(a\Omega \subset \Omega_{\mathrm{B}}\)，这蕴含 \(a \in \mathfrak{c}\)。于是映射 \(a \mapsto a\omega\) 诱导出从 \(\mathrm{A}/\mathfrak{c}\) 到 \(\Omega/\Omega_{\mathrm{B}}\) 的一个单射；从而有 \(\mathrm{long}(\mathrm{A}/\mathfrak{c}) \leqslant \mathrm{long}(\Omega/\Omega_{\mathrm{B}}) = \mathrm{long}(\mathrm{B}/\mathrm{A})\)。

若 long(A/c) = long(B/A)，则 Aω + Ω＿B = Ω。可设理想 c 含于 m＿A（否则 A 等于 B，从而是 Gorenstein 环）。由于 Ω＿B = cΩ 含于 m＿AΩ，由 Nakayama 引理，ω 生成 Ω。于是 A-模 Ω 是循环模，从而是秩为 1 的自由模，这意味着 A 是 Gorenstein 环（第 4 节，命题 7 的推论 1）。

\begin{enumerate}
\def\labelenumi{\Alph{enumi})}
\setcounter{enumi}{3}
\tightlist
\item
  现在处理一般情形。用 A′ 表示环 A{]}X{[}，即（IX, App., 第 2 节）多项式环 Λ{[}X{]} 在素理想 m\_AΛ{[}X{]} 处的局部环；它是一维的平坦整 A-代数，其剩余域 κ\_A′ 等同于 κ\_A(X)，而分式域等同于 K(X)（前引文）。由第 3 节命题 5 的推论，A′-模 A′ ⊗\_A Ω 是对偶化的。令 B′ = A′ ⊗\_A B；它是 Λ′ 在 K(X) 中的整闭包（V, § 1, 第 3 节，命题 13 及第 5 节，命题 16）。传递子 c′ = A′ : B′ 等于 cA′（I, § 2, 第 10 节，公式 (11)）。对每个有限长度的 A-模 M，有 long\_A′(A′ ⊗\_A M) = long\_A(M)：事实上，由于 A-代数 A′ 是平坦的，只需对 M 是单模（即同构于 κ\_A）的情形证明此关系；但在这种情形，A′ ⊗\_A κ\_A 等同于 κ\_A′，断言得证。于是我们有
\end{enumerate}

\[
\operatorname{long} _ {\mathrm{A}} (\mathrm{B} / \mathfrak {c}) = \operatorname{long} _ {\mathrm{A} ^ {\prime}} \left(\mathrm{B} ^ {\prime} / \mathfrak {c} ^ {\prime}\right) \quad \text { and } \quad \operatorname{long} _ {\mathrm{A}} (\mathrm{B} / \mathrm{A}) = \operatorname{long} _ {\mathrm{A} ^ {\prime}} \left(\mathrm{B} ^ {\prime} / \mathrm{A} ^ {\prime}\right).
\]

环 A′ 满足命题的假设，且它的剩余域是无限的。由证明的 C) 部分，有 \(\text{long}_{\text{A}'}(\text{B}'/\text{c}') \leqslant 2 \text{long}_{\text{A}'}(\text{B}'/\text{A}')\)，且等号成立蕴含 A′ 是 Gorenstein 环；而后者又蕴含 A 是 Gorenstein 环（§ 3, 第 8 节，命题 12 的推论 1）。

